\section{Harmonic models, density, and the actual cloud maps}
\label{sec:harmonic-model}

\status{HM1--HM12 identify the literal polynomial lattices, prove the
bounded harmonic comparison and its density consequence, and construct
the complementary-middle cloud maps and coefficient-ring transport.
The HH rank input is RU11; its numerical supply and the final applications
of the cloud maps are separate. No integral semisimplicity is assumed.}

\paragraph{HM1--HM2: an explicit harmonic quotient.}
Work in the source coefficient setting $w^2=p$, with $v_p(w)=1/2$.
Put $k=m+1$, $S_d=\binom{d+k-1}{k-1}$ and $S_{-1}=0$.
For a $k\times k$ matrix $Z$, let
$c_i(Z)=(-1)^i\det Z_{\widehat i,\{0,\ldots,k-2\}}$.
The actual two-letter map sends $x_i$ to the first-column entry
of one matrix and $y_i$ to $c_i$ of an independent matrix.
Its leading monomials recover every exponent of $x,y$, so it is
injective and saturated over $O$. Identifying the matrices gives
\[
 \beta_{ab}:P_{ab}\longrightarrow H_{(a,0,\ldots,-b)},\qquad
 \ker\beta_{ab}=hP_{a-1,b-1},\quad h=\sum_i x_i y_i,
 \qquad H_{ab}=S_aS_b-S_{a-1}S_{b-1}.                    \tag{HM.1}
\]
Indeed CB6 gives surjectivity, the repeated-column determinant gives
$\beta(h)=0$, and CB3's dimension formula gives equality of kernels.
The cofactor identity
$c(g^{-1}Z)=\det(g)^{-1}g^t c(Z)$ checks the actual determinant twists
for every lattice $gO^k$.

Let $\Delta=\sum_i\partial_{x_i}\partial_{y_i}$.
The commutator $[\Delta,h]F=(a+b+k)F$ gives
\[
 \pi_{AB}=\sum_{\ell=0}^{\min(A,B)}
 \frac{(-1)^\ell}{\binom{A+B+k-2}{\ell}}
 \frac{h^\ell}{\ell!}\frac{\Delta^\ell}{\ell!},\qquad
 E_j=\binom{A+B+k-j-1}{j}^{-1}
 \frac{h^j}{j!}\pi_{A-j,B-j}\frac{\Delta^j}{j!}.        \tag{HM.2}
\]
These project onto $h^j\ker\Delta$ and sum to the identity.
Both divided operators preserve the lattice with basis
$x^\alpha y^\beta/\alpha!$. Legendre's formula gives
$v_p\binom nr\le\lfloor\log_p n\rfloor$; comparison with the
ordinary coefficient lattice costs at most $(k-1)\lfloor\log_p(A+B+k)\rfloor$.
Thus both actual maps $\operatorname{pr}_j,\operatorname{io}_j$
obtained from (HM.2) have denominator at most
\[
 U=\left\lceil(m+2)\bigl(\log_p(8a+m)+1\bigr)\right\rceil,
 \qquad \sum_j\operatorname{io}_j\operatorname{pr}_j=1,
                                                               \tag{HM.3}
\]
when $A+B+1\le8a$. Division by $h^j$ costs nothing: it has a
monic leading term. Equivariance gives the same bound at every lattice.

\paragraph{HM3--HM5: normalization and rank splitting.}
Tensor over $N=|TT|$ disjoint letters and put $u=NU$.
Write $A_H=\mathrm{AHH}$, $A_j=\mathrm{Ar}(\lambda_j)$,
$d_j=A_H-A_j\ge0$. For cleaned vectors $x,\Phi$ the literal maps are
\[
 x_j=p^u w^{d_j}\operatorname{pr}_jx,\qquad
 \Phi_j=p^{2u+d_j}\Phi\circ(w^{-d_j}\operatorname{io}_j),\qquad
 \mathrm{be}_j(x_j,\Phi_j)
   =p^{3u}\mathrm{be}_H(\operatorname{io}_j\operatorname{pr}_jx,\Phi).
                                                               \tag{HM.4}
\]
Applying the lattice bounds to the source, midpoint, target and every
competing edge verifies all cleaned conditions, with unchanged height
and conductor.

If integral matrices satisfy $\sum_jB_j=zI_n$, the total number of
their Smith factors with valuation at most $v_p(z)$ is at least $n$.
Otherwise their small-factor coordinate rows have a common primitive
kernel vector, contradicting $\sum_jB_jx=zx$.
Consequently an HH certificate at depth $D$ splits among the actual
harmonic components at depth $D+3u\le D+4u$.

There are $N_e=(p^{n'k}-1)/(p^{n'}-1)$ outgoing rank-one edges.
Transport by $\det(g)^{-\lambda_{\rm last}}\operatorname{rowAct}(g^{-1})$
preserves the actual lattices and pairing. A certificate at one edge,
with primal family in $\mathrm{CCL}$ and dual family merely in $R'_e$,
therefore gives one at every edge. If $2D<c$, dividing each diagonal
block by its scalar makes the joint pairing matrix the identity modulo
$w$: different-edge pairings lie in $w^cO$. Hence
\[
                         n\le\dim V(\lambda)/N_e.       \tag{HM.5}
\]
This proves the one-sided \texttt{nLeD} as well as its cleaned-dual version.

\paragraph{HM6--HM7: the actual density inequality.}
Use the fixed-weight numerical bounds supplied by NP6, set
$D_1=D_0+4NU$, and apply RU11's HH mass
\[
                         (1-k_x(A,B))\dim HH/N_e.
\]
For $J=\lfloor J_0\rfloor$, decreasing all degrees by at most $J$
lowers the segment-minimum height by at most $NJ_0$.
Thus $t_*(\alpha,\beta)=\lceil t_0(\alpha,\beta)+NJ_0\rceil$
admits every small-shift certificate. Combining (HM.5), the defining
infimum $k_F$, and the telescoping identity
$\sum_j\prod_sH_{A_s-j_s,B_s-j_s}=\prod_sS_{A_s}S_{B_s}$ gives
\[
 \sum_{0\le j_s\le J} k_F(A-j,B-j)
       \prod_s H_{A_s-j_s,B_s-j_s}
 \le k_x(A,B)\prod_sS_{A_s}S_{B_s}.                    \tag{HM.6}
\]
For $R_j=S_\alpha S_\beta/(S_{\alpha+j}S_{\beta+j})$,
monotonicity of $t/(t+k-1)$ gives
\[
 \sum_{j=0}^{J}\frac{H_{\alpha\beta}}
                         {S_{\alpha+j}S_{\beta+j}}
 \ge1-R_{J+1}
 \ge1-\exp\!\left(-\frac{(k-1)(J+1)}{\beta+J+k}\right).\tag{HM.7}
\]
Summing (HM.6) under $(\alpha,\beta,j)\mapsto(\alpha+j,\beta+j,j)$
therefore proves the source \texttt{Avg}: if $\delta$ bounds (HM.7)
uniformly and the extended product box has size at most $2\delta^N$
times the original box, then $\operatorname{mean}k_F\le2\operatorname{mean}k_x$.
RD11--RD15 supply these numerical inequalities.

\paragraph{HM8: actual duality and disjoint tensor factors.}
The transported row-homogeneous CB2 basis $x_i$ has target exponents
$\pi_i$ and orders $o_i=A-\pi_i\ge0$. Thus
\[
 R_e=\bigoplus_iOw^{-o_i}x_i,\qquad
 R'_e=\bigoplus_iOw^{-o_i}p^{-\pi_i}x_i^*,\qquad
 p^Aw^{-2o_i}p^{-\pi_i}=1.                            \tag{HM.8}
\]
The order inequality counts entries in the first edge rows of each
tableau. This proves the exact dual-height and perfect-pairing interfaces.
Splitting letter blocks into $TT$ and its complement gives actual
\texttt{RestTensor}: multiply the disjoint polynomial factors, multiply
their basis evaluations on the dual side, and multiply basis cardinalities.
Transported product bases prove both lattice equalities, including midpoints.

\paragraph{HM9--HM11: complementary pairing and cloud maps.}
For complementary middle shapes $\alpha'_i=M-\alpha_{r-1-i}$, put
$e=\min\alpha$, $\gamma=\alpha-e$, $T=M-e$ and let $J_r$ reverse rows.
Jacobi's complementary-minor identity proves that
\[
 \tau f(Y)=\det(Y)^T f(Y^{-t}J_r)                      \tag{HM.9}
\]
is an integral isomorphism from shape $\alpha'$ to $\gamma$.
IW3--IW5 bound the normalized Fock dual lattice by
$p^\kappa M_\gamma\subset M_\gamma^\#\subset M_\gamma$,
$\kappa=\binom r2\lfloor\log_p(\gamma_0+r-1)\rfloor$.
Take $\kappa=0$ for the one-dimensional case.
Smith reduction of $p^\kappa B_0$, followed by rounding its valuations
upward, gives integral families with pairings $\delta_{ab}p^{\kappa_b}$,
$0\le\kappa_b\le\kappa$.
The pairing $L(f,f')=p^\kappa B_0(f/\det^e,\tau f')$
has simultaneous row covariance $\det(g)^M$.

The actual contraction is coefficientwise
\[
 \Lambda(v,t)(Z)=L_{Y,Y'}\bigl(v(Z\operatorname{diag}(1,Y,1))
                 t(Z\operatorname{diag}(1,Y',1))\bigr).\tag{HM.10}
\]
Lift each middle tableau minor by adjoining row and column zero,
and add the required power of $Z_{00}$. The resulting integral
$s_b,s'_a$ have row-zero degree $2c$ and last-row degree zero.
The highest-line calculation, evaluated at $Z=I$, gives exactly
$\Lambda(Fs_b,s'_a)=\delta_{ab}p^{\kappa_b}H_\zeta F$.
Tensoring produces the source cloud maps with $\#b=\dim\mathrm{Cl}$.
Multiplication by the primitive $H_\zeta$ has exact integral preimage,
by Gauss valuation. The actual dual map is
$g\mapsto p^{Nc}g\circ L_a$; testing against other target edges
costs precisely $2Nc$ powers of $w$.
This proves full \texttt{RestCloudMaps}, with conductor $c_2-2Nc$.

\paragraph{HM12: coefficient-ring independence.}
Hensel-lift a common residue generator. The powers
\[
                          1,\theta,\ldots,\theta^{n'-1}
\]
form a $\mathbb Z_p$-basis of each unramified $W$.
A common quadratic residue lift $\eta$ gives the basis
$1,\eta,w,\eta w$ of $O$ over $W$; successive reduction modulo $w$
proves independence, and completeness proves spanning.
Matching these bases preserves their defining equations and sends $w$ to $w$.
Uniqueness of the Hensel lifts makes the isomorphism commute with Frobenius.
Coefficient transport then preserves every actual polynomial lattice,
dual, normalized pairing and integral rank certificate. Thus the literal
\texttt{LPprime} statement is independent of the chosen RS setting.
